\BourbakiExerciseGroup

\begin{enumerate}
\def\labelenumi{\arabic{enumi})}
\tightlist
\item
  A point \(a \in \mathbb{R}\) is said to be left adherent to a subset \(\mathbf{A}\) of \(\mathbb{R}\) if it lies in the closure of \(\mathbf{A}\) and if there is an interval \(]a, b[ (a < b)\) not containing any point of \(\mathbf{A}\) . Show that the set of left adherent points of a subset \(\mathbf{A}\) of \(\mathbb{R}\) is countable (set up a one-to-one correspondence between the set of left adherent points and a set of pairwise disjoint open intervals). Hence show that every well-ordered subset of \(\mathbb{R}\) is countable.
\end{enumerate}

¶ 2) A countable dense subset \(\mathbf{A}\) of \(\mathbb{R}\) is not closed. {[}If \((a_{n})\) is a sequence obtained by arranging the points of \(\mathbf{A}\) in some order, define a sequence of intervals \([b_{n}, c_{n}]\) such that \(b_{n-1} < b_{n} < c_{n} < c_{n-1}\) for all \(n\) , and such that \([b_{n}, c_{n}]\) contains no point \(a_{k}\) with index \(k \leqslant n\) ; then use Theorem 2.{]} Deduce that \(\mathbb{R}\) is uncountable (cf.~§ 8, no. 6).

\begin{enumerate}
\def\labelenumi{\arabic{enumi})}
\setcounter{enumi}{2}
\item
  Let \((\mathbf{I}_n)\) be an infinite sequence of non-empty open intervals in \(\mathbf{R}\) such that \(\overline{\mathbf{I}}_n \cap \overline{\mathbf{I}}_m = \emptyset\) for \(m \neq n\) . Show that the complement of the union of the \(\mathbf{I}_n\) is a perfect set (Chapter I, § 2, no. 6); in particular, Cantor's triadic set is perfect.
\item
  The sum of the lengths of the intervals contiguous to Cantor's triadic set is equal to 1. Define a totally disconnected perfect subset of {[}0, 1{]} such that the sum of the lengths of the intervals contiguous to A and contained in {[}0, 1{]} is a given number m such that \(0 < m \leqslant 1\) .
\item
  Let \(l\) be a number \(> 0\) , and suppose that to each point \(x \in \mathbf{R}\) corresponds an open interval \(\mathbf{I}(x)\) with centre \(x\) and length \(\leqslant l\) . Show that every compact interval \([a, b]\) can be covered by a finite number of intervals \(\mathbf{I}(x_i)\) , the sum of whose lengths is \(\leqslant l + 2(b - a)\) . (Prove that if the proposition is true for each interval \([a, x]\) such that \(a \leqslant x < c\) , then there exists \(d > c\) such that it is true for each interval \([a, y]\) such that \(a \leqslant y < d\) .) If \(l = (b - a)/n\) , where \(n\) is an integer \(\geqslant 1\) , show that the result cannot be improved.
\end{enumerate}

¶ 6) a) Let X be a non-empty linearly ordered set, endowed with the topology \(\mathcal{C}_{0}(X)\) (Chapter I, § 2, Exercise 5). Show that X is compact if and only if each subset of X has a least upper bound, in other words if and only if X is a complete lattice (Set Theory, Chapter III, § 1, Exercise 11). (To show that the condition is necessary, argue as for Theorem 3; for sufficiency, consider a filter \(\mathfrak{F}\) on X and show that, if A is the set of greatest lower bounds of sets of \(\mathfrak{F}\) , then the least upper bound of A is a cluster point of \(\mathfrak{F}\) ).

\begin{enumerate}
\def\labelenumi{\alph{enumi})}
\setcounter{enumi}{1}
\tightlist
\item
  Give an example of a complete lattice X which is not compact in the topology \(\mathcal{C}_{0}(\mathbf{X})\) .
\end{enumerate}

¶ 7) Let X be a non-empty linearly ordered set, endowed with the topology \(\mathfrak{G}_{\mathfrak{0}}(\mathbf{X})\) .

\begin{enumerate}
\def\labelenumi{\alph{enumi})}
\tightlist
\item
  If X is connected, show that it has the following two properties:
\end{enumerate}

α) Every non-empty subset of X which is bounded above has a least upper bound (if A is non-empty and bounded above, and if B is the set of upper bounds of A, and C the set of lower bounds of B, show that \(B \cup C = X\) and that B and C are closed).

β) The set X is without gaps, in other words (Set Theory, Chapter III, § 1, Exercise 18) every open interval {]}a, b{[} (a \textless{} b) of X is non-empty (use the argument of Theorem 4).

\begin{enumerate}
\def\labelenumi{\alph{enumi})}
\setcounter{enumi}{1}
\item
  Show conversely that if X has properties \(\alpha\) and \(\beta\) ), then X is connected. (Show first that every closed interval \([a, b]\) is connected: assuming that this interval is the disjoint union of two non-empty closed sets A and B, and that \(a \in A\) , consider the greatest lower bound of B and hence obtain a contradiction).
\item
  Show that if X is connected it is locally compact and locally connected, and that the only connected subsets of X are the intervals (bounded or unbounded).
\end{enumerate}

\begin{enumerate}
\def\labelenumi{\arabic{enumi})}
\setcounter{enumi}{7}
\tightlist
\item
  Let X and Y be two linearly ordered sets, endowed with the topologies \(\mathcal{G}_{0}(X)\) and \(\mathcal{G}_{0}(Y)\) respectively.
\end{enumerate}

\begin{enumerate}
\def\labelenumi{\alph{enumi})}
\tightlist
\item
  Suppose \(\mathbf{X}\) connected and let \(f: \mathbf{X} \to \mathbf{Y}\) be a continuous mapping. Show that, for all \(x, y\) in \(\mathbf{X}\) such that \(x < y\) , every \(z \in \mathbf{Y}\) belonging to the closed interval with end-points \(f(x)\) and \(f(y)\) belongs to \(f(\mathbf{X})\)
\end{enumerate}

(use Exercise 7). Deduce that \(f\) is a homeomorphism of \(\mathbf{X}\) onto \(f(\mathbf{X})\) if and only if \(f\) is continuous and strictly monotone.

\begin{enumerate}
\def\labelenumi{\alph{enumi})}
\setcounter{enumi}{1}
\tightlist
\item
  Give an example of a homeomorphism of the rational line Q onto itself which is not strictly monotone.
\end{enumerate}

\begin{enumerate}
\def\labelenumi{\arabic{enumi})}
\setcounter{enumi}{8}
\tightlist
\item
  \begin{enumerate}
  \def\labelenumii{\alph{enumii})}
  \tightlist
  \item
    Let X be a countable linearly ordered set without gaps (Exercise 7). Show that there is a strictly increasing bijective mapping \(\varphi\) of X onto one of the four intervals of Q with end-points o and i. {[}Arrange the elements of X and the points of the appropriate interval of Q in sequences \((a_n), (b_n)\) , and define \(\varphi\) by induction.{]} If X is endowed with the topology \(\tilde{C}_0(X)\) , then \(\varphi\) is a homeomorphism of X onto \(\varphi(X)\) .
  \end{enumerate}
\end{enumerate}

\begin{enumerate}
\def\labelenumi{\alph{enumi})}
\setcounter{enumi}{1}
\item
  Deduce from a) that every countable dense subset of an open interval {]}a, b{[} of R is homeomorphic to Q.
\item
  Deduce from a) that if X is any countable ordered set endowed with the topology \(\mathfrak{G}_{0}(\mathbf{X})\) , then there exists a strictly increasing homeomorphism of X onto a subspace of Q {[}embed X in a countable linearly ordered space \(X'\) without gaps, such that \(\mathfrak{G}_{0}(\mathbf{X})\) is induced by \(\mathfrak{G}_{0}(\mathbf{X}')\) {]}.
\end{enumerate}

\begin{enumerate}
\def\labelenumi{\arabic{enumi})}
\setcounter{enumi}{9}
\tightlist
\item
  Show that every countable subset of Cantor's triadic set K, which is dense in K and contains no end-point of an interval contiguous to K, is homeomorphic to the rational line Q (use Exercise 9).
\end{enumerate}

¶ 11) a) Let X be a linearly ordered set with the topology \(\mathfrak{G}_{0}(\mathbf{X})\) . If X is connected and has a countable dense subset A, show that there exists a homeomorphism (strictly monotonic by virtue of Exercise 8) of X onto one of the intervals of R with end-points o, i, which maps A onto the intersection of Q with this interval (use Exercises 9 and 7).

\begin{enumerate}
\def\labelenumi{\alph{enumi})}
\setcounter{enumi}{1}
\item
  Deduce from a) that if B is a subset of R whose complement is dense in R, then there exists a homeomorphism of R onto itself which maps B onto a subset of the set CQ of irrational numbers (consider a dense countable subset A of R contained in CB).
\item
  In the set \(\mathbf{R} \times \{0, 1\}\) , linearly ordered by the lexicographic ordering, let \(\mathbf{R}'\) denote the complement of \(\mathbf{Q} \times \{1\}\) . Show that \(\mathbf{R}'\) , endowed with the topology \(\mathfrak{G}_0(\mathbf{R}')\) , is locally compact and totally disconnected, and contains a countable dense subset \(\mathbf{Q}'\) homeomorphic to \(\mathbf{Q}\) , but that the topology induced on a non-empty open interval of \(\mathbf{R}'\) does not have a countable base; in particular such an interval cannot be homeomorphic to any subset of \(\mathbf{R}\) .
\end{enumerate}

¶ 12) a) Let A be a well-ordered set and let I denote the interval {[}0, 1{[} of R. Show that the set X = A × I, linearly ordered by the lexicographic ordering and endowed with the topology \(\mathfrak{G}_{0}(X)\) , is connected (Exercise 7); there are therefore linearly ordered sets X which are connected in the topology \(\mathfrak{G}_{0}(\mathbf{X})\) and have an arbitrary cardinal {[}cf.~§ 4, Exercise 7 b){]}.

\begin{enumerate}
\def\labelenumi{\alph{enumi})}
\setcounter{enumi}{1}
\tightlist
\item
  Take A to be an uncountable well-ordered set in which every segment \(\leftarrow, t]\) is countable: the corresponding topological space X is called the Alexandroff half-line. Show that for every interval \([a, b]\) of X (a \textless{} b) there exists a strictly increasing homeomorphism of \([a, b]\) onto the interval \([o, i]\) of R. (Prove by transfinite induction that there exists a strictly increasing homeomorphism of \(A \cap [a, b]\) , A being identified with the set of points \((t, o)\) of X onto a subspace of \([o, i]\) .)
\end{enumerate}

¶ 13) Let \(a\) be an irrational number \(>0\) . For each rational number \(x\) , let \(f_{a}(x)\) be the real number in the interval \([0, a]\) such that \(x - f_{a}(x)\) is an integral multiple of \(a\) . Show that \(f_{a}\) is a continuous injective mapping of \(\mathbf{Q}\) into \([0, a]\) . Deduce with the help of Exercise 9 that there exist continuous bijective mappings of \(\mathbf{Q}\) onto itself, whose inverse mapping is not continuous at any point.

¶ 14) Let X be a connected Hausdorff topological space not consisting of a single point, and let \(\Delta\) denote the diagonal of \(X \times X\) .

\begin{enumerate}
\def\labelenumi{\alph{enumi})}
\item
  Show that there is no partition of \(\mathbb{C}\Delta\) into two open sets \(\mathbf{D}_1, \mathbf{D}_2\) such that \(\mathbf{D}_1^{-1} = \mathbf{D}_1\) and \(\mathbf{D}_2^{-1} = \mathbf{D}_2\) . {[}Begin by remarking that \(\overline{\mathbf{D}_1(x)}\) and \(\overline{\mathbf{D}_2(x)}\) are connected for each \(x \in \mathbf{X}\) (use Exercise 4 a) of Chapter I, § 11); deduce that, if \(y \in \mathbf{D}_1(x)\) , we have \(\overline{\mathbf{D}_2(x)} \times \{y\} \subset \mathbf{D}_1\) , and show that this implies that, for each \(x \in \mathbf{X}\) , one or the other of the sets \(\mathbf{D}_1(x)\) , \(\mathbf{D}_2(x)\) is empty; hence show that one or the other of the sets \(\mathbf{D}_1, \mathbf{D}_2\) is empty, which is contrary to hypothesis.{]} Deduce that either \(\mathbb{C}\Delta\) is connected or else \(\mathbb{C}\Delta\) has exactly two components A, B such that \(\mathbf{B} = \overline{\mathbf{A}}\) .
\item
  A linear ordering on X is said to be compatible with the topology \(\mathfrak{C}\) of X if \(\mathfrak{C}_0(\mathbf{X})\) is coarser than \(\mathfrak{C}\) . Deduce from \(a\) ) that there exists such a linear ordering on X if and only if \(\mathbb{C}\Delta\) is not connected, and that there are then exactly two such linear orderings. {[}With the notation of \(a\) ), show that the order relation must be either \((x, y) \in \mathbf{A}\) or \((x, y) \in \mathbf{B} = \overline{\mathbf{A}}\) .{]} Show that the intervals, for these orderings, are connected in the topology \(\mathfrak{C}\) , and that they are the only connected subsets of X with respect to \(\mathfrak{C}\) .
\item
  Show that if X is either locally connected or locally compact in the topology \(\mathfrak{C}\) , and if there exists a linear ordering on X compatible with \(\mathfrak{C}\) , then we must have \(\mathfrak{C} = \mathfrak{C}_0(X)\) . {[}If X is locally connected, use b). If X is locally compact, show that every compact neighbourhood of a point \(x \in X\) is also a neighbourhood of \(x\) with respect to \(\mathfrak{C}_0(X)\) , by considering the frontier of the neighbourhood{]}.
\end{enumerate}

* d) Let X be the subspace of \(\mathbf{R}^2\) consisting of (o, o) and all pairs \((x, y)\) such that \(x > 0\) and \(y = \sin (\mathrm{i} / x)\) . Show that there exists a linear ordering on X which is compatible with the topology induced on X by the topology of \(\mathbf{R}^2\) , but that this topology is strictly finer than \(\mathcal{C}_0(X)\) . *

Give an example of a connected Hausdorff space X, no point of which has a fundamental system of connected neighbourhoods, and on which there is a linear ordering compatible with the topology of X {[}cf.~Chapter I, § 11, Exercise 2 b){]}.

¶ 15) a) Let X be a connected Hausdorff space such that, for each \(x \in X\) , \(\complement\{x\}\) has exactly two components. Show that, if K is either of these components, \(K = K \cup \{x\}\) {[}Chapter I, § 11, Exercise 4 a){]}. Let \(x, y\) be two distincts points of X, let A and B be the components of \(\complement\{x\}\) and let \(A', B'\) be the components of \(\complement\{y\}\) . Show that one of the two sets A, B is contained in \(A'\) or in \(B'\) .

\begin{enumerate}
\def\labelenumi{\alph{enumi})}
\setcounter{enumi}{1}
\tightlist
\item
  Let \(x_0\) be a point of \(X\) and let \(A(x_0), B(x_0)\) be the components of \(\complement\{x_0\}\) . For each \(x \neq x_0\) there is exactly one of the two components of \(\complement\{x\}\) which is either contained in \(A(x_0)\) or contains \(A(x_0)\) ; denote this component by \(A(x)\) . Show that the relation \(A(x) \subset A(y)\) is a linear order relation compatible (Exercise 14) with the topology of \(X\) . c) Extend the conclusion of b) to the case where \(\complement\{x\}\) has two components for all \(x \in X\) except for one point \(a \in X\) , for which \(\complement\{a\}\) is connected.
\end{enumerate}

* \(d\) ) Let \(\mathbf{X}\) be the subspace of \(\mathbf{R}^2\) consisting of the points \(a = (0, 1)\) , \(b = (0, -1)\) and all pairs \((x, y)\) such that \(x \neq 0\) and \(y = \sin(1/x)\) . Show that \(\mathbf{X}\) is connected in the topology \(\mathfrak{G}\) induced on \(\mathbf{X}\) by the topology of \(\mathbf{R}^2\) , that \(\mathbb{G}\{a\}\) and \(\mathbb{G}\{b\}\) are connected and that, for each point \(x\) of \(\mathbf{X}\) other than \(a\) and \(b\) , \(\mathbb{G}\{x\}\) has exactly two components; but that there exists no linear ordering on \(\mathbf{X}\) compatible with the topology \(\mathfrak{G}_{*}\) .

¶ 16) Let X be a connected Hausdorff space which is completely irreducible between two distinct points \(a\) and \(b\) , that is to say such that there is no connected subset of X other than X itself which contains both \(a\) and \(b\) . Show that \(\mathbf{X}' = \mathbb{C}\{a, b\}\) is connected and that, for each \(x \in \mathbf{X}'\) , the complement of \(x\) in X has exactly two components \(\mathbf{A}(x)\) and \(\mathbf{B}(x)\) such that \(a \in \mathbf{A}(x)\) , \(b \in \mathbf{B}(x)\) , \(a \notin \mathbf{B}(x)\) , \(b \notin \mathbf{A}(x)\) . {[}Show that \(\mathbb{C}\{x\}\) cannot be partitioned into three non-empty open sets, by using Exercise 4 a) of Chapter I, § 11.{]} Deduce that there exists a linear ordering on X which is compatible with the topology of X (use Exercise 15).

¶ 17) Let X be a connected Hausdorff space such that whenever X is covered by three non-empty connected sets distinct from X, there are two of these sets whose union is distinct from X.

\begin{enumerate}
\def\labelenumi{\alph{enumi})}
\item
  Show that, for every point \(x \in \mathbf{X}\) , \(\mathbb{G}\{x\}\) has at most two components (same method as in Exercise 16).
\item
  Show that there are at most two points of X whose complements are connected. Moreover, if there are two distinct points a, b with this property, show that, for each x other than a or b, a and b belong to different components of \(C\{x\}\) .
\item
  Deduce from a) and b) that there is a linear ordering on X compatible with the topology of X (use Exercise 15).
\end{enumerate}

¶ 18) a) Let X be a compact, connected, locally connected space which is irreducible between two of its points \(a\) , \(b\) (Chapter II, § 4, Exercise 19). Let \(x\) be a point of X other than \(a\) or \(b\) . Show that \(\mathbb{C}\{x\}\) has exactly two components and that \(a\) and \(b\) belong to different components of \(\mathbb{C}\{x\}\) . {[}Arguing as in Exercise 16, show first that \(\mathbb{C}\{x\}\) cannot have more than two components. Next, if V is any connected closed neighbourhood of \(x\) which contains neither \(a\) nor \(b\) , let \(\mathrm{A}_{\mathrm{V}}\) , \(\mathrm{B}_{\mathrm{V}}\) be the components of \(a\) , \(b\) respectively in \(\mathbb{C}\mathrm{V}\) ; show that \(\mathrm{A}_{\mathrm{V}} \neq \mathrm{B}_{\mathrm{V}}\) by using Exercise 17 of Chapter II, § 4; finally consider the union of the sets \(\mathrm{A}_{\mathrm{V}}\) (resp. \(\mathrm{B}_{\mathrm{V}}\) ) as V runs through the set of all connected closed neighbourhoods of \(x\) in X.{]} Deduce that there exists a linear ordering on X such that \(\mathfrak{C}_{0}(\mathbf{X})\) is the given topology on X {[}Exercises 15 and 14 c){]}.

\begin{enumerate}
\def\labelenumi{\alph{enumi})}
\setcounter{enumi}{1}
\tightlist
\item
  Let X be a compact connected space which is irreducible between two of its points a, b. Show that if X has more than one prime constituent (Chapter II, § 4, Exercise 20) then the space \(X'\) of prime constituents of X (loc. cit.) is irreducible between the prime constituents of a and b, and that there is a linear ordering on \(X'\) such that the topology of \(X'\) is identical with \(\mathcal{G}_{0}(\mathbf{X}')\) . * Give an example in which \(X'\) is distinct from X {[}cf.~Exercise 15 d{]}. *
\end{enumerate}

¶ 19) a) Let X be a connected Hausdorff space and let \(a, b\) be two distinct points of X such that there is no connected closed subset of X, other than X itself, which contains both \(a\) and \(b\) . If \(x\) is a point of X other than \(a\) or \(b\) , then \(\mathbb{C}\{x\}\) is not connected if and only if, for each \(y \in \mathbb{C}\{x\}\) , there is in \(\mathbb{C}\{y\}\) a connected subset which is closed in X and contains \(x\) and one of the points \(a, b\) . {[}To show that the condition is necessary, use Exercise 4 a) of Chapter I, § 11; to show that it is sufficient, let A (resp. B) be the set of all \(y \in X\) such that there is a connected subset of \(\mathbb{C}\{y\}\) which is closed in X and contains \(a\) and \(x\) (resp. \(b\) and \(x\) ); show that A and B are non-empty and open and disjoint from each other, and that \(\mathbb{C}\{x\} = A \cup B\) .

\begin{enumerate}
\def\labelenumi{\alph{enumi})}
\setcounter{enumi}{1}
\tightlist
\item
  Let X be a compact connected space which is irreducible between two of its points \(a, b\) (Chapter II, § 4, Exercise 19). Show that if for each pair of distinct points \(x, y\) of X there is a compact connected set in X which contains \(y\) and one of \(a, b\) but does not contain \(x\) , then there is a linear ordering on X such that \(\mathfrak{C}_{0}(\mathbf{X})\) is the given topology on X (Exercises 15 and 16).
\end{enumerate}

\begin{enumerate}
\def\labelenumi{\arabic{enumi})}
\setcounter{enumi}{19}
\item
  In the construction of Exercise 23 of Chapter I, § 9, take \(\mathbf{X}_0\) to be the interval {[}0, 1{]} of R. Show that the homeomorphisms of X onto itself are the same as the homeomorphisms of \(\mathbf{X}_0\) onto itself, although the topology of X is strictly finer than the topology of \(\mathbf{X}_0\) {[}use Exercise 20 b) of Chapter I, § 8{]}.
\item
  Show that if \(r > 2\) there is no \(r\) -fold transitive group of homeomorphisms of \(\mathbf{R}\) onto itself (consider the subgroup leaving two points of \(\mathbf{R}\) fixed).
\end{enumerate}

